\addcontentsline{toc}{chapter}{Conclusions}
\chapter*{Conclusions}

The detailed analysis and implementation work conducted in this project have confirmed that while the digital manufacturing industry has advanced significantly in terms of hardware, the software infrastructure for machine control remains rooted in the mid-20th century. The reliance on G-code creates a significant barrier to efficiency, customization, and safety.

The Problem-Based Learning (PBL) investigation highlighted three critical issues: the lack of readability and modifiability of machine instructions, the absence of abstraction which hides design intent, and the steep learning curve required to debug or optimize manufacturing processes manually. The proposed solution—the GG-Code DSL—emerges as a necessary and viable innovation.

By introducing high-level programming concepts such as variables, loops, functions, and semantic operations into the manufacturing workflow, the DSL transforms machine instructions from cryptic data streams into readable, maintainable software code. The successful specification of the formal grammar (Chapter 2) and the implementation of a robust parser (Chapter 3) demonstrate the technical feasibility of this approach. The parser is now capable of interpreting complex motion commands, geometric primitives, and algorithmic logic, validating the core architectural decisions.

In conclusion, the project has progressed from problem definition to a functional prototype of the language frontend. The grammar and parser provide a solid foundation for the next phase: the development of the semantic analyzer and the code generator backend, which will translate the validated Abstract Syntax Tree into optimized G-code for target machines. This tool has the potential to significantly streamline the prototyping workflow and empower users to leverage the full capabilities of their manufacturing hardware.
